% =============================================================================
% lp_workflow_body.tex
% TikZ body for the LP-FLARE workflow figure.
% Two fully independent side-by-side columns:
%   LEFT  -- Constraint-First Screening (LP, this work)
%   RIGHT -- Composition-First Screening (conventional brute-force)
% Each column has its own Inputs box and its own Survivor box; no shared boxes
% and no arrows crossing between the two chains.
% =============================================================================

\begin{tikzpicture}[
    >={Stealth[length=1.8mm, width=1.5mm]},
    font=\sffamily,
    node distance = 4mm and 12mm,
    inbox/.style = {
        rectangle, draw, rounded corners=1.5pt, align=center,
        text width=33mm, minimum height=15mm, inner sep=2pt,
        font=\scriptsize, fill=gray!8
    },
    bfbox/.style = {
        rectangle, draw, rounded corners=1.5pt, align=center,
        text width=33mm, minimum height=10mm, inner sep=2pt,
        font=\scriptsize, fill=red!5
    },
    bfexp/.style = {
        rectangle, draw=red!75, line width=0.8pt, rounded corners=1.5pt, align=center,
        text width=33mm, minimum height=12mm, inner sep=2pt,
        font=\scriptsize, fill=red!15
    },
    lpbox/.style = {
        rectangle, draw, rounded corners=1.5pt, align=center,
        text width=33mm, minimum height=10mm, inner sep=2pt,
        font=\scriptsize, fill=blue!5
    },
    outbox/.style = {
        rectangle, draw, rounded corners=1.5pt, align=center,
        text width=33mm, minimum height=12mm, inner sep=2pt,
        font=\scriptsize, fill=green!8, line width=0.7pt
    },
    arrow/.style = {->, semithick, line width=0.5pt},
    coltitle/.style = {font=\footnotesize\bfseries, align=center}
]

% =========================================================================
% LEFT column: Constraint-First Screening (LP, this work)
% =========================================================================

\node[coltitle] (lp-title) {Constraint-First\\Screening};

\node[inbox, below=1.5mm of lp-title] (lp-in) {%
    \textbf{Inputs}\\[1pt]
    Ti-V-Ta-Nb-Mo-Zr-\\
    Cr-Hf-Fe-Re-W\\
    design space};

\node[lpbox, below=of lp-in] (lp1) {%
    \textbf{Build linear constraints}\\
    \((\mathbf{G},\mathbf{h})\),\;\(\sum_i x_i = 1\)};

\node[lpbox, below=of lp1] (lp2) {%
    \textbf{Feasible polytope}\\
    \(P=\{\mathbf{x}\in\Delta:\mathbf{G}\mathbf{x}\le\mathbf{h}\}\)};

\node[lpbox, below=of lp2] (lp3) {%
    \textbf{Enumerate feasible alloys}\\
    5\,at.\% grid compositions\\
    contained in \(P\)};

\node[outbox, below=of lp3] (lp-out) {%
    \textbf{Survivor compositions}\\[1pt]
    61{,}507 compositions};

\draw[arrow] (lp-in) -- (lp1);
\draw[arrow] (lp1)   -- (lp2);
\draw[arrow] (lp2)   -- (lp3);
\draw[arrow] (lp3)   -- (lp-out);

% =========================================================================
% RIGHT column: Composition-First Screening (conventional brute-force)
% =========================================================================

\node[coltitle, right=of lp-title] (bf-title) {Composition-First\\Screening};

\node[inbox, below=1.5mm of bf-title] (bf-in) {%
    \textbf{Inputs}\\[1pt]
    Ti-V-Ta-Nb-Mo-Zr-\\
    Cr-Hf-Fe-Re-W\\
    design space};

\node[bfbox, below=of bf-in] (bf1) {%
    \textbf{Generate compositions}\\
    (5\,at.\% resolution)\\
    \(\sim 3{\times}10^{7}\) alloys};

\node[bfexp, below=of bf1] (bf2) {%
    \textbf{Evaluate FLARE}\\ for every composition\\[1pt]
    \textcolor{red!70!black}{$\sim 3\times10^{7}$ model calls}};

\node[bfbox, below=of bf2] (bf3) {%
    \textbf{Apply filter}\\
    compare each response\\
    against thresholds};

\node[outbox, below=of bf3] (bf-out) {%
    \textbf{Survivor compositions}\\[1pt]
    61{,}507 compositions};

\draw[arrow] (bf-in) -- (bf1);
\draw[arrow] (bf1)   -- (bf2);
\draw[arrow] (bf2)   -- (bf3);
\draw[arrow] (bf3)   -- (bf-out);

\end{tikzpicture}
